\chapter{Kritische Reflexion und Ausblick}
Kommen wir nun zur kritischen Reflexion und zum Ausblick.

\section{Kritische Reflexion}
In dieser Arbeit wurde ein Machine-Learning-Modell entwickelt und trainiert, um zwischen den Aktivitäten Gehen, Joggen und Rennen zu unterscheiden. Die Ergebnisse zeigen, dass das Modell in der Lage ist, diese Aktivitäten mit hoher Genauigkeit zu klassifizieren. Besonders hervorzuheben ist die hohe Genauigkeit bei der Klassifikation der Aktivität Gehen, während die Klassifikation von Rennen noch Verbesserungspotenzial aufweist.

Die Analyse der Feature-Wichtigkeit hat gezeigt, dass bestimmte Merkmale, wie die Standardabweichung und Extremwerte der Beschleunigung entlang der x-Achse, besonders wichtig für die Klassifikation sind. Dies deutet darauf hin, dass die Bewegungsdynamik in Längsrichtung ein entscheidender Faktor für die Unterscheidung der Aktivitäten ist. Die Erkenntnisse aus dieser Analyse können in zukünftigen Arbeiten genutzt werden, um die Modellkomplexität zu reduzieren und die Genauigkeit weiter zu verbessern.

Ein kritischer Punkt dieser Arbeit ist die Qualität und Repräsentativität der Trainingsdaten. Obwohl die Daten sorgfältig gesammelt und vorbereitet wurden, könnte eine Erweiterung des Datensatzes, insbesondere für die Klasse Rennen, die Klassifikationsgenauigkeit weiter verbessern. Zudem könnte die Berücksichtigung weiterer sensorischer Merkmale, wie z.B. Herzfrequenz oder Atemfrequenz, zusätzliche Informationen liefern und die Klassifikation weiter verbessern.

\section{Ausblick}
In zukünftigen Arbeiten könnte das entwickelte Modell in ein Echtzeitsystem integriert werden, das kontinuierlich Sensordaten aufnimmt und die aktuelle Aktivität klassifiziert. Dies würde es ermöglichen, die körperliche Aktivität in Echtzeit zu überwachen und Feedback zu geben. Ein solches System könnte in verschiedenen Anwendungsbereichen eingesetzt werden, wie z.B. im Sport, in der Rehabilitation oder im Gesundheitswesen.

Ein weiterer wichtiger Schritt wäre die Implementierung einer Funktion, die es ermöglicht, ein Datenpaket aufzunehmen und an das Modell zu senden, um die aktuelle Aktivität zu klassifizieren. Dies würde die praktische Anwendbarkeit des Modells erheblich erhöhen und den Weg für die Entwicklung von Echtzeit-Anwendungen ebnen.

Zusammenfassend bietet diese Arbeit eine solide Grundlage für die Klassifikation von körperlichen Aktivitäten mittels Machine Learning. Die gewonnenen Erkenntnisse und entwickelten Modelle können in zukünftigen Arbeiten weiter verbessert und in praktischen Anwendungen eingesetzt werden. Die Integration des Modells in ein Echtzeitsystem und die Implementierung einer Funktion zur Klassifikation von Datenpaketen sind vielversprechende nächste Schritte, die die praktische Anwendbarkeit und den Nutzen dieser Arbeit weiter erhöhen würden.
